\documentclass[sigconf,review,anonymous]{acmart}
\providecommand{\tightlist}{\setlength{\itemsep}{0pt}\setlength{\parskip}{0pt}}
\setcopyright{none}
\acmConference[MMSys '27]{ACM Multimedia Systems Conference}{2027}{}
\emergencystretch=3em
\sloppy
\begin{document}
\title{Structural Construction of Retrieval Graphs for Long-Video Question Answering}
\author{Anonymous Author(s)}
\affiliation{\institution{Anonymous Institution}\country{}}
\begin{abstract}
Graph-based pipelines for long-video question answering increasingly build their retrieval
graph with a large multimodal model, paying a generative cost per chunk that grows with
video duration. We build the graph structurally instead --- from compressed-domain signals and
CLIP embeddings, on CPU, with no generative model in the construction path --- and verify the
construction rather than benchmark it. A block-diagonal construction that enumerates only
within-scene node pairs is bit-identical to dense-then-prune construction (same nodes,
edges, weights at zero tolerance, and PageRank) while visiting 0.2\% of the pairs, making
construction \textasciitilde234\ensuremath{\times} faster and \textasciitilde6.8\ensuremath{\times} smaller in peak memory at 4,892 nodes, on a
complete-block edge configuration (our accuracy numbers come from other configurations). Query latency
separates with graph size --- exponent 2.06 (95\% CI {[}2.02, 2.13{]}) dense versus 0.83
({[}0.80, 0.88{]}) sparse --- and past a certain size dense retrieval fails to return within
3600 s. We also report where this does not help: a captioning stage dominates query time,
so a three-orders-of-magnitude retrieval speedup yields no end-to-end gain, and accuracy
sits in the weakly-supervised band (Acc@GQA 0.167 on a held-out validation split, \textasciitilde3.4B
answerer on CPU). Two falsified hypotheses --- codec-derived scene boundaries and codec-based
frame admission, each tested against a stated kill criterion --- show that the block
structure, not the codec signal, carries the result.
\end{abstract}
\ccsdesc[500]{Information systems~Video search}
\ccsdesc[300]{Computing methodologies~Video summarization}
\keywords{long-video question answering, retrieval graphs, compressed-domain video, efficiency}
\maketitle
\section{Introduction}

Answering questions about long video requires structure. An hour of footage yields tens of
thousands of candidate frames, and flat retrieval over that pool becomes intractable before
it becomes inaccurate. The standard response is to build an intermediate graph over frames,
shots, or scenes and retrieve against it.

Increasingly, that graph is written by a large multimodal model. EgoSG prompts a proprietary
model once per fixed-duration chunk to emit a symbolic scene description; Vgent extracts
entities per clip with an LVLM and links clips that share them. Both produce richer
representations than ours, but both make construction a model-metered operation whose cost
scales with duration, and the symbolic line places a proprietary dependency at the base of
the pipeline that its own authors report does not transfer to open-source substitutes. For a
surveillance archive, or any footage that cannot leave the premises, that is a poor
foundation.

Two obvious claims are not ours. That construction cost is query-independent and amortises
across questions is Vgent's argument, measured before us; that compressed-domain signals
can stand in for full decode is older still (§2). Our contribution is narrower. We build the
graph with no generative model in the construction path --- a CLIP image encoder over
admitted frames is the only network --- and we verify the construction rather than benchmark
it.

\textbf{Construction.} The scene-sparse graph keeps only within-scene edges, so the conventional
construction --- materialise the dense graph, then prune across scenes --- spends almost all its
work on edges it discards: at N=4,892 admitted frames over 528 scenes it visits 11.96 million pairs
to keep 23,571. Building the diagonal blocks directly visits only the pairs it keeps, and
yields the same graph rather than an approximation of it: identical nodes and edges, zero
edge-attribute mismatches at tolerance 0.0, bit-identical PageRank, and identical
personalised-PageRank output across five seeds, with zero behavioural differences
over 526 grounded-QA questions. Construction time falls from 49.71 s to 0.212 s
(\textasciitilde234\ensuremath{\times}) and peak memory from 6.33 GB to 0.93 GB (\textasciitilde6.8\ensuremath{\times}). This is measured on the
complete-block edge configuration; no accuracy number we report is measured on it (§8).

\textbf{Scaling.} The same block structure changes how query cost grows. Over 31 clips under real
text queries, dense retrieval latency scales with graph size at exponent 2.063 (95\% CI
{[}2.020, 2.125{]}) and scene-sparse at 0.834 ({[}0.801, 0.884{]}). Past a certain size
this becomes a tractability boundary: under a 3600 s / 29 GB watchdog, dense retrieval never
returned on clips of 6,559 and 13,506 admitted frames, while sparse retrieval answered them in
0.010 s and 0.021 s.

\textbf{What does not improve.} We state two limits here because they are results, not
concessions. First, the retrieval speedup does not reach the user. At N=4,892, retrieval
mechanics fall from 8.535 s to 0.030 s, but a captioning-and-verification stage of 44--67 s,
O(top\_k) rather than O(N), dominates the query; end to end, the sparse arm is slower at the
median (0.802\ensuremath{\times}). Much recent work reports retrieval-side speedups without measuring the
pipeline around them. We measure it, and we do not project a crossover, because our own
data contradict the assumption such a projection needs (§5.3). Second, accuracy sits in the
weakly-supervised band: held-out Acc@GQA on NExT-GQA is 0.1667 {[}0.088, 0.243{]} with a \textasciitilde3.4B
answerer on CPU. Our claim is that the pipeline answers within that band, not that
it advances it.

\textbf{What does not matter.} We set kill criteria for two hypotheses we expected to confirm,
and both triggered. Codec-derived scene boundaries do not beat content-blind ones at matched
segment count (pre-registered; +0.088 M-Avg, 95\% CI {[}\ensuremath{-}0.009, +0.204{]} on long videos), and
codec-based frame admission is indistinguishable from uniform sampling at matched budget
(criterion recorded, though we cannot show it preceded the run). Neither where the
video is cut nor which frames are kept does the work; the block structure does. A pipeline
can therefore adopt the construction without our codec signal, segmentation, or admission
policy.

\textbf{Contributions.}

\begin{enumerate}
\def\labelenumi{\arabic{enumi}.}
\tightlist
\item
  A block-diagonal graph construction shown bit-identical to dense-then-prune, at \textasciitilde234\ensuremath{\times}
  lower build time and \textasciitilde6.8\ensuremath{\times} lower peak memory on the complete-block configuration (§4).
\item
  Query-latency scaling exponents with bootstrap confidence intervals over 31 clips, and a
  tractability boundary where dense retrieval does not complete (§5.1--§5.2).
\item
  A construction path with no generative model: frame selection performs zero neural
  forward passes, verified by instrumentation; full ingest costs 11.976 s per
  video, of which 2.845 s is selection (§3).
\item
  An Amdahl accounting of the query path, including a caption-cache locality cost that our
  sparse retrieval pays (§5.3--§5.4).
\item
  Two negative results under stated kill criteria, delimiting which components are
  load-bearing (§7).
\end{enumerate}

\section{Related Work}

\textbf{The token bottleneck.} Multimodal LLMs are bounded by input token capacity, so applying
them directly to long video forces aggressive subsampling: at 1 FPS, published per-question
frame budgets run from 32 frames (InternVL3) and 180 (VideoLLaMA3) to 512 (Qwen2.5-VL) and
3,900 (Gemini Flash 2.0), against the tens of thousands of frames in an hour of footage.
Every approach is in some form a strategy for deciding what to discard: subsample and accept
the loss, retrieve per query, or build an intermediate representation that later queries
reuse. We sit in the third family and borrow from the second at query time.

\subsection{Graph-structured intermediate representations}

\textbf{EgoSG} {[}Taluzzi et al., 2026{]} partitions video into fixed 60 s chunks at 1 FPS and prompts
Gemini Flash 2.0 to emit a symbolic scene graph per chunk --- objects, spatial relations,
timestamped action hyperedges --- each chunk's graph produced by prompting the model to update
the previous one. On HD-EPIC VQA it improves over raw-video input by 3.22 points for Gemini
itself and 5.60 for Qwen2.5-VL-14B.

EgoSG's graph is a symbolic reasoning substrate and ours an embedding retrieval structure;
relational queries spanning a whole video genuinely favour the former, so we compare on
construction cost, not representational power.

Three properties of that construction matter. First, it is \textbf{model-metered and sequential}:
roughly 5.7 s per one-minute clip against a proprietary API, scaling linearly with duration,
and because chunk \(i\)'s graph depends on chunk \(i-1\)'s it cannot be parallelised across a
long video. The authors further report that open-source substitutes struggle to produce
accurate graphs reliably, and name that dependency as a limitation. Our scenes are
independent, so the block structure is parallel by definition, and no model participates in
construction.

Second, \textbf{no correctness guarantee is available for the artifact}: a manual audit of five
clips finds roughly 5\% of nodes and relations in error, 15.5\% false action hyperedges and
13.7\% of objects missing, and the authors note that optimising generation quality is not
their focus. That is the expected behaviour of a generative construction, not a criticism. A
structural construction can instead be verified exactly, and is (§4.1) --- though only as a
\emph{fidelity} guarantee that our block-diagonal path reproduces our own dense-then-prune
reference, not as semantic correctness. It cannot hallucinate an object, but neither can it
recognise one.

\textbf{Vgent} {[}Shen et al., NeurIPS 2025{]} is the closer foil. It partitions video into 64-frame
clips, extracts entities per clip with an LVLM, and links clips sharing merged prototype
entities. It rests on the same amortisation argument we make --- construction is offline and
query-independent, so one build serves every question asked of that video --- and, unusually,
reports the split explicitly: 20.13 s of query-independent construction and 3.93 s of
query-dependent work per minute of video, against 20.81 s per minute for the query-dependent
Video-RAG baseline, a claimed 1.73\ensuremath{\times} end-to-end speedup on VideoMME. We are therefore not the
first to argue that construction cost deserves its own measurement, and do not claim to be.

Where we differ is what construction costs. Vgent invokes an LVLM per clip as EgoSG invokes
Gemini per chunk; our frame selection runs on CPU with zero neural forward passes, and our
only construction-time network is CLIP ViT-B/32 over admitted frames --- an image encoder, not
a generative model (§3). \textbf{We draw no cost ratio against any of these systems.} Our
selection stage costs 2.845 s per video and full ingest 11.976 s, while Vgent's 20.13 s is
per \emph{minute of video} and includes LVLM extraction; the denominators do not match, so the
honest claim is about the \emph{class} of model each construction requires, not about seconds.

Three further systems --- \textbf{EGAgent} {[}Rege et al., 2026{]}, \textbf{MemDreamer} {[}Chen et al., 2026{]}
and \textbf{GraphVideoAgent} {[}Chu et al., ACM MM 2025{]} --- build their graphs with an LLM extractor,
a Gemini-3.1-Pro extractor, or by parsing a captioning model's output, so each places a
generative pass inside construction or one step upstream of it. None reports a
construction-cost figure, so no ratio comparison is available; the claim is about model
class only.

\subsection{Grounded video QA, and the accuracy frontier}

\textbf{NExT-GQA} {[}Xiao et al., 2023{]} requires that a system answer correctly \emph{and} point at its
evidence, reporting Acc@GQA alongside mIoP, IoP@0.5 and mIoU. Its finding that strong
answerers often ground poorly --- right for the wrong reasons --- is why we report grounded
accuracy rather than answer accuracy alone (§6). The weakly-supervised band there is occupied
by Temp{[}CLIP{]} with NG+ (16.0 Acc@GQA), SeViLA as reproduced by the benchmark authors (16.6),
LangRepo (17.1) and FrozenBiLM with NG+ (17.5). Published figures use the full 5,553-question
test set, where ours is a 120-question held-out carve of the validation split (§6).

A more recent line attaches an explicit temporal grounder, often with cooperating agents:
\textbf{MUPA} reports 28.7 Acc@GQA at 2B and 30.3 at 7B, \textbf{VideoMind} 25.2 at 2B. These define
the current frontier. We do not compete with them on accuracy --- our contribution is
construction cost and our answer quality sits a band below (§6) --- and cite them to situate
that gap rather than to select a weaker comparison set.

\subsection{Compressed-domain video analysis}

Using bitstream signals --- motion vectors, residual energy, packet size --- to avoid full decode
is long-standing, running from motion-vector surrogates for optical flow {[}Zhang et al., CVPR
2016{]} through CoViAR and DMC-Net. It has recently reached video language models:
\textbf{CoPE-VideoLM} {[}Sarkar et al., 2026{]} converts each P-frame into eight compact \ensuremath{\Delta}-tokens from
its motion vectors and residuals, reporting up to 86\% lower time-to-first-token and up to 93\%
fewer visual tokens across 14 benchmarks.

Our use of the codec differs in kind, and that distinction is what makes our negative result
compatible with their positive one. CoPE-VideoLM uses codec primitives as a \emph{representation},
encoded and consumed by the model, at the cost of a trained \ensuremath{\Delta}-encoder and end-to-end
fine-tuning. We use codec statistics as a \emph{selection and weighting heuristic} over frames we
then treat conventionally, training nothing. That codec-derived salience does not beat
uniform sampling for \emph{admission} (§7.2) says nothing about whether the same fields carry
signal when \emph{encoded}; we take no position on theirs.

Our claim is correspondingly modest, and §7 constrains it: compressed-domain signals here
are a \emph{cheap} source of structure, not an \emph{informative} one.

\subsection{Positioning}

Against EgoSG and its family we replace a frontier-model-generated symbolic graph with a
structurally-constructed one: built with no generative model, verifiable, parallel, locally
runnable --- and semantically poorer. Against grounded-QA baselines we land in the weakly-supervised band
while training nothing and running on CPU; against the agentic frontier we are a band behind
on accuracy. The claim is narrow: for
retrieval over long video, the intermediate graph does not have to be expensive, and does not
have to be generated.

\section{Method}

IRIS ingests a video once into a compact index, then answers arbitrary text queries against
it. Throughout, \textbf{N} denotes the number of \emph{surviving} frames after admission, not the frame
count in the container --- the distinction matters for every scaling figure we report.

\subsection{Ingest and frame admission}

Ingest demuxes the video and computes per frame a scalar \textbf{action score} combining a codec
packet-size residual, motion energy and luma entropy at weights 0.5 / 0.3 / 0.2. Packet size
stands in for the residual signal because it is what the decoder API exposes: libavcodec
exposes motion vectors as side data (\texttt{EXPORT\_MVS}) but no coefficient-level residual,
per-block quantisation parameter or reconstruction-error signal, so a frame-level
packet-size proxy is the only residual-shaped channel this tooling admits. Frames are
assigned to tiers by thresholding that signal (\texttt{salient\_thresh} 0.35, \texttt{candidate\_thresh}
0.08, adaptive per-video thresholding enabled), and local maxima are marked as peaks.
Admitted frames --- the survivors --- are the input to everything downstream; production
retention across the annotated corpus falls between 10.41\% and 11.14\%.

\textbf{No neural network runs during frame selection.} We verify this by instrumentation rather
than inspection: monkeypatching \texttt{torch.nn.Module.\_\_call\_\_} for the duration of every
\texttt{parse\_video()} call and counting invocations across 32 videos yields exactly zero.

CLIP embeddings (ViT-B/32) are computed for admitted frames \emph{after} that stage, and are what
the zero-forward-pass claim excludes. Captioning is deferred entirely; \texttt{caption} remains
\texttt{None} at ingest. Over the same frozen 32-video sample:

\begin{table}[t]
\footnotesize
\setlength{\tabcolsep}{4pt}
\caption{Ingest stage wall time over the frozen 32-video sample.}
\begin{tabular}{lrrr}
\toprule
stage & mean wall & min & max \\
\midrule
\texttt{parse\_video} (selection) & 2.845 s & 0.314 s & 14.870 s \\
CLIP enrichment & 9.131 s & 0.501 s & 56.885 s \\
\textbf{full ingest} & \textbf{11.976 s} & 0.814 s & 71.754 s \\
\bottomrule
\end{tabular}
\end{table}

CLIP enrichment, at 26.56 ms per admitted frame, costs roughly three times the selection stage.
\textbf{The honest headline for construction is therefore \textasciitilde12 s per video, not \textasciitilde2.9 s.} Our ingest
as a whole is \emph{not} forward-pass-free; the smaller figure describes the forward-pass-free
portion only, and any comparison against a system whose reported construction cost includes
its neural extraction stage is not like-for-like on it. Selection-stage peak RSS is 530 MB
mean, measured in a run where torch was never loaded.

\subsection{Scene segmentation}

Survivors are partitioned into scenes. The default rule derives boundaries from valleys in
the packet-size curve taken directly from the container --- no pixel decode, no model call.
Segmentation is a configuration axis with four settings (\texttt{codec}, \texttt{fixed\_count},
\texttt{fixed\_seconds}, \texttt{fixed\_time\_matched}), which are exactly the four arms compared in §7.1. We
present the codec rule as one cheap option, not a superior one: §7.1 finds boundary placement
within noise across all four at matched segment count. What the pipeline depends on is that
survivors are partitioned at all, not that the partition is content-adaptive.

\subsection{The scene-sparse graph and block-diagonal construction}

Let S scenes partition the N survivors, scene \(i\) holding \(n_i\). The scene-sparse graph
connects every pair within a scene and no pair across scenes, giving an adjacency that is
block-diagonal under a scene-ordered permutation. Edges carry semantic, motion and temporal
components. The edge count is

\[|E| = \sum_i n_i(n_i-1)/2\]

against N(N\ensuremath{-}1)/2 for the dense graph. The conventional construction computes the dense graph
and prunes cross-scene edges; ours enumerates the diagonal blocks directly. At N=4,892 over
S=528 scenes those quantities are 23,571 and 11,963,386: the pruning construction discards
\textbf{99.8\%} of the pairs it examines, while block-diagonal construction visits exactly as many
pairs as it retains.

In the balanced case, with n = N/S per scene, the pair count is S\ensuremath{\cdot}n(n\ensuremath{-}1)/2 \ensuremath{\approx} \(N^2/(2S)\) against
\(N^2/2\), so the saving is the scene count S. Two boundaries follow: a single scene (S=1) makes
the paths identical and saves nothing, and for fixed N and S the saving is maximal when
scenes are equal in size, since \(\sum_i n_i^2\) is minimised at \(n_i = N/S\) --- an uneven partition with
one dominant scene approaches the dense cost. At N=4,892 over S=528 the realised ratio is
0.197\% against 1/S = 0.189\%, so this partition is close to balanced.

Because the pruned graph and the block-diagonal graph are the same object, this is a change
of construction order, not of representation; §4.1 verifies that empirically rather than
resting on the argument. This section describes the \textbf{complete-block} graph: the production
edge mode (\texttt{hierarchical\_sparse}, §3.5) shares the block structure but populates each block
sparsely, so \textbar E\textbar{} above is an upper bound on it rather than a description of it.

\subsection{Retrieval}

A text query is encoded with the same CLIP model used at ingest and resolved in stages:

\begin{enumerate}
\def\labelenumi{\arabic{enumi}.}
\tightlist
\item
  \textbf{Scene shortlist.} Each scene is represented by the centroid of its frame embeddings;
  only the top \texttt{w} scenes by query--centroid score proceed, with \(w = \max(4, \lceil\sqrt{S}\rceil)\).
\item
  \textbf{Descent.} Within shortlisted scenes, frames are ranked by personalised PageRank over
  the scene-sparse graph, seeded from the query, with damping 0.5 and a rank-space blend
  (\ensuremath{\lambda} = 0.5) between semantic and codec-derived rank.
\item
  \textbf{Shortcut.} Scenes are ranked by best per-frame CLIP similarity; if the margin between
  the top two shortlisted scenes exceeds \ensuremath{\tau} (\texttt{scene\_shortcut\_margin}, default 0.015), the
  anchor scene's exact top-k is returned and the union-subgraph PPR is skipped. When only
  one scene is shortlisted the margin is set to infinity, so the shortcut fires
  deterministically.
\item
  \textbf{Top-k, captioning, answering.} Retrieved frames are captioned and passed to the
  answerer (\texttt{granite4:micro}, \textasciitilde3.4B, Q4\_K\_M, CPU).
\end{enumerate}

Two properties are load-bearing later. The \textbf{shortlist} makes retrieval sublinear by
construction --- it is why §5.1's exponent is below 1 --- but it is also a hard recall gate: a
scene excluded here cannot contribute a frame downstream (§8). The \textbf{shortcut} bypasses the
graph entirely and is the mechanism behind the guard violations in §5.1; both of its trigger
conditions become more likely as S falls, consistent with the two smallest clips in the
corpus being the two that trip it. It is exact rather than approximate --- it returns the same top-k the
full path would --- so no retrieved result is degraded by it.

\subsection{Configuration and provenance}

Results are produced under three settings of \texttt{graph\_edge\_mode}, and we state per experiment
which was used rather than describing a single frozen configuration:

\begin{table*}[t]
\small
\setlength{\tabcolsep}{4pt}
\caption{Graph configuration used by each experiment.}
\begin{tabular}{lll}
\toprule
experiment & \texttt{graph\_mode} & \texttt{graph\_edge\_mode} \\
\midrule
§4.1 identity gate, §4.2 build savings & \texttt{scene\_sparse} & \texttt{fully\_connected} \ensuremath{\leftrightarrow} \texttt{block\_diagonal} \\
§5.1--5.2 scaling, tractability & \texttt{flat} / \texttt{scene\_sparse} & \texttt{fully\_connected} \\
§5.3--5.4 end-to-end, caption stage & \texttt{scene\_sparse} & \texttt{hierarchical\_sparse} \\
§6.2 MLVU, §7.1 segmentation ablation & \texttt{scene\_sparse} & \texttt{hierarchical\_sparse} \\
§6.1 NExT-GQA grounded QA & \texttt{flat} & --- \\
\bottomrule
\end{tabular}
\end{table*}

All three edge modes produce within-scene-only graphs under \texttt{graph\_mode="scene\_sparse"}:
\texttt{fully\_connected} and \texttt{block\_diagonal} connect every intra-scene pair, while
\texttt{hierarchical\_sparse} builds a tiered edge set (temporal at window 1, hierarchy parents,
salient-semantic top-4, motion-neighbour top-2) and then removes every cross-scene edge. The
block structure is common to all three.

\textbf{They do not produce the same graph within the blocks, and two consequences follow that we
state rather than gloss.} On the same clip, \texttt{fully\_connected} fills 23,571 intra-scene edges
where \texttt{hierarchical\_sparse} fills 6,284, roughly 3.75\ensuremath{\times} sparser. First, \textbf{§4's construction
saving is measured on the complete-block configuration, which produced no accuracy number in
this paper} --- §6.2 and §7.1 run \texttt{hierarchical\_sparse} and §6.1 a flat dense graph (§4.3
scopes this). Second,
\textbf{\texttt{hierarchical\_sparse} pays its own quadratic pass} that block-diagonal construction does
not remove: motion-neighbour selection scans all N(N\ensuremath{-}1)/2 pairs to keep two per node, and the
salient-semantic pass scans all pairs among salient nodes. The tiered mode is sparse in its
\emph{output}, not in its \emph{construction} (§8).

\textbf{The scaling argument and the accuracy results therefore rest on different graph
constructions, and that gap is open.} Every cached index behind §5.1 is built under
\texttt{fully\_connected} or \texttt{block\_diagonal}; not one contains a \texttt{temporal}, \texttt{hierarchy\_*},
\texttt{semantic\_salient} or \texttt{motion\_neighbor} edge. The degree bound §5.1 fits transfers to the
tiered path only through the shared block structure, not within a scene. Under
\texttt{graph\_mode="flat"} the block-diagonal path is unreachable, since the flat construction never
partitions nodes into groups; §6.1's grounded-QA results are a dense-graph measurement,
reported as a correctness floor rather than as a measurement of the sparse construction.
Per-experiment configuration hashes recorded in the artifacts, not this table, are the
reproducibility anchor.

\subsection{Edge weights}

Every edge carries three components: semantic, a ReLU-clamped cosine between CLIP embeddings;
motion, a normalised action-score gap \(\max(0, 1 - |a_u - a_v|/R)\) over the per-build range
R; and temporal, \(1/(1 + |t_u - t_v|)\) in seconds. Each edge family combines them with fixed
coefficients, except \texttt{fully\_connected}, which uses \ensuremath{\alpha}\ensuremath{\cdot}semantic + \ensuremath{\beta}\ensuremath{\cdot}motion (\ensuremath{\alpha}=0.4, \ensuremath{\beta}=0.3 in
every cached index); the per-family formulas are listed with the released artifacts. One
consequence for §5.1: \ensuremath{\alpha} and \ensuremath{\beta} affect the dense arm's edge weights and not the scene-sparse
arm's at all.

\section{Construction Cost}

The scene-sparse graph is the same graph the conventional construction produces, built at a
fraction of the cost. We establish the identity first and report the savings second: a
construction speedup is uninteresting if it silently changes what is built.

One scope statement belongs here rather than in a footnote. This section concerns the
\textbf{complete-block} configuration, in which every intra-scene pair carries an edge. No accuracy
number in this paper is measured on it: §6.2 and §7.1 use the tiered \texttt{hierarchical\_sparse}
configuration, which shares the block structure but populates each block sparsely (6,284
edges against 23,571 on the same clip), and §6.1 a flat dense graph (§3.5). §4.3 states what follows from that.

\subsection{The identity gate}

Two paths produce a scene-sparse graph: the reference path materialises the dense N\ensuremath{\times}N graph
and prunes cross-scene edges; the block-diagonal path materialises only the within-scene
blocks. At N=4,892 admitted frames over 528 scenes (VIRAT) they agree exactly: edge
count 23,571 on both paths, matching the theoretical prediction; identical node sets,
with zero edges present in only one graph; zero mismatches on every edge attribute
(\texttt{weight}, \texttt{semantic\_weight}, \texttt{motion\_weight}, \texttt{temporal\_weight}, \texttt{edge\_type}) at tolerance
\textbf{0.0}, not at a floating-point epsilon; bit-identical PageRank across all 4,892 nodes; and
identical personalised-PageRank top-20 ordering and scores across five seeds.

The two paths are also indistinguishable at the end of the pipeline. Over 526 NExT-GQA
validation questions they produce bit-identical retrieval --- same retrieved order, peak
frame, predicted span, IoP and peak-in-gold on every question (peak-in-gold 0.3175, mIoP
0.3140 under both) --- and both index caches were sha256-identical before and after the
run. We run this second gate because graph identity does not by itself guarantee identical
downstream behaviour.

What the gate establishes is an equivalence between two construction paths, not a re-run of
a committed result: no committed grounding number was produced under the complete-block
configuration, since none of our accuracy arms uses it (§3.5).

This is an identity rather than an approximation because cross-scene edges are absent from
the scene-sparse graph by definition. The reference path computes them and then discards
them; the block-diagonal path never computes them.

\subsection{Savings}

\begin{table}[t]
\footnotesize
\setlength{\tabcolsep}{4pt}
\caption{Construction cost at $N=4{,}892$, $S=528$ (complete-block configuration).}
\begin{tabular}{lrrr}
\toprule
 & reference & block-diag. & ratio \\
\midrule
wall-clock (per build) & 49.71 s & 0.212 s & \textbf{234.4\ensuremath{\times}} \\
peak RSS & 6.33 GB & 0.93 GB & \textbf{6.80\ensuremath{\times}} \\
\bottomrule
\end{tabular}
\end{table}

The two rows come from two protocols, because wall-clock and peak memory are not well
measured the same way.

\textbf{Wall-clock.} 159 builds on one clip (N=4,892, 528 scenes), arms alternated across three
rounds so that machine drift affects both equally, with load and free memory logged before
each arm. Round-level ratios were 235.21\ensuremath{\times}, 233.99\ensuremath{\times} and 234.45\ensuremath{\times}, pooled median 234.45\ensuremath{\times} ---
agreement within 0.52\%. The library code under test is pinned to a commit; the harness and its
driver were committed afterwards, so the measurement is reproducible from the
repository but was not executed from a pinned harness. An earlier execution of the identical
protocol, on an uncommitted tree, pooled 233.53\ensuremath{\times} --- agreement within 0.4\%, a cross-session
stability point stronger than either run's internal agreement.

\textbf{Peak RSS.} Five repeats per arm, each a single build in a fresh process, medians
reported: 6,325,764,096 bytes for the reference path and 930,693,120 bytes for the
block-diagonal path. A peak RSS taken over 50 consecutive in-process builds is a peak over
the loop, not over a build. Unlike the wall-clock ratio, this figure is not tied to a
commit. An independently recorded post-dedup ratio (6.8039\ensuremath{\times}) agrees to three significant
figures, and every build in both runs produced exactly 23,571 edges.

\textbf{Three caveats travel with the table; the first makes it conservative.} First, the fast
arm is timed over 50 consecutive in-process builds. This puts the timed quantity safely
above timer resolution but charges each iteration roughly 0.09 s of fixed allocation
overhead --- negligible against a 50 s build, substantial against a 0.2 s one. A single cold
build in a fresh process, which is what production performs, measures faster (\ensuremath{\approx}0.12 s). We
report the looped figure because it is the reproducible one, and the effect is to understate
the ratio. This does leave a looped timing beside a cold-process memory figure; each
protocol was chosen for the validity of its own quantity, and we have not measured which
direction a looped memory figure would move the ratio.

Second, this isolates the \texttt{\_build\_graph} stage with frames and embeddings already cached, so
it is not a full-ingest figure. The selection stage that produces those inputs costs 2.845 s
mean wall time at 530 MB mean peak RSS over a 32-video stratified sample, CPU-only, with
zero neural forward passes; full ingest including CLIP enrichment costs
11.976 s (§3). A full-ingest ratio would therefore be far smaller than the build ratio; we do
not compute one, because that sample and this clip differ.

Third, the reported ratio is a within-session figure. The reference arm's wall time has
varied across sessions running the same corrected code --- 61.39 s and 48.57 s in two earlier
fresh-process runs, against 49.71 s here, a spread of roughly 20\% that the interleaved
protocol neither explains nor reproduces, with machine state logged idle throughout
(4.0--11.9\% CPU, 18.6--19.3 GB free of 33.5 GB). Interleaving establishes that the two arms
drift together within a session; it does not establish that the reference arm's absolute
cost is stable across sessions, and the 0.52\% agreement should be read as within-session
precision.

Two corrections belong on the record. An earlier version of this measurement reported \textasciitilde220\ensuremath{\times}
and \textasciitilde6.7\ensuremath{\times}; both paths then performed a redundant edge-and-PageRank pass whose result was
immediately discarded, and removing it changed the timing of both arms without changing what
they build --- verified bit-identically on nodes, all five edge fields, per-node \texttt{scene\_id} and
PageRank, and independently by §4.1. Intermediate point measurements of the corrected code
ranged from 259.23\ensuremath{\times} (single fresh-process probe) to a 396.73\ensuremath{\times} five-repeat median, a spread
dominated by the fast arm's cold-build noise (0.1203 s to 0.2164 s, \textasciitilde1.8\ensuremath{\times} on a
sub-quarter-second operation, against under 1\% on the reference arm). That spread is why we
report the looped, interleaved protocol rather than any point measurement.

\subsection{What the construction claim covers}

Within the configuration it is measured on, the construction saving cannot be absorbed by a
downstream stage, because there is no downstream stage inside construction. That is the
sense in which it is the more durable of our two efficiency results: §5.3 reports a
retrieval speedup of three orders of magnitude that nonetheless fails to reach the user.

The limit, stated plainly. The saving is measured on the complete-block configuration, and
no reported accuracy number comes from it: MLVU and the segmentation ablation use the tiered
configuration, and NExT-GQA a flat dense graph to which block-diagonal construction cannot
apply at all. The tiered configuration shares the block structure but not the edge set, and
block-diagonal construction does not make it cheaper to build: that mode selects neighbours by scanning all pairs and then keeping a few,
so it pays a quadratic pass of its own (§3). Carrying the result into the tiered
configuration means bounding those scans to the scene partition --- a change the current
implementation does not make, and one that would produce a \emph{different} graph rather than an
identically-constructed one, since global top-k followed by cross-scene pruning is not the
same as within-scene top-k. It therefore requires accuracy validation rather than an
identity gate, and we mark it as future work (§8).

\section{Query Scaling}

\subsection{Latency exponents}

We fit log(latency) = k\ensuremath{\cdot}log(N) + c per arm over 31 UCF-Crime clips under \textbf{real CLIP text
queries} (50 per clip), with clip-level bootstrap confidence intervals (seed 42, B=10,000,
resampling clips within each N-bucket and refitting).

\begin{table}[t]
\footnotesize
\setlength{\tabcolsep}{4pt}
\caption{Query-latency scaling exponents (31 clips, real CLIP text queries).}
\begin{tabular}{llrrlr}
\toprule
arm & fit range & n points & k & 95\% CI & \(R^2\) \\
\midrule
dense (flat) & N \ensuremath{\geq} 1102 (headline) & 4 & 2.063 & {[}2.020, 2.125{]} & 0.9994 \\
scene-sparse & N \ensuremath{\geq} 1102 (headline) & 6 & 0.834 & {[}0.801, 0.884{]} & 0.9830 \\
dense (flat) & full range & 17 & 1.980 & {[}1.951, 2.006{]} & 0.9978 \\
scene-sparse & full range & 19 & 0.497 & {[}0.461, 0.530{]} & 0.8980 \\
\bottomrule
\end{tabular}
\end{table}

Over the measured range, dense query cost grows with the square of graph size and
scene-sparse cost grows sublinearly. The headline intervals are disjoint and the
scene-sparse upper bound lies below 1.0, the pre-registered bar for that separation.

We report two scene-sparse exponents and take as headline the one \emph{less} favourable to us.
The full-range fit (0.497) gives a larger separation than we report. We infer from the fit's
shape alone --- shallower small-N slope, poorer explanatory power (\(R^2\) 0.898 against 0.983) ---
that small-N clips sit near a noise floor where fixed per-query overhead dominates; we have
not verified this directly, and no repeated small-N runs or fitted overhead constant exist.
The N \ensuremath{\geq} 1102 fit is both the regime this paper is about and the conservative choice.

\textbf{Four caveats travel with the table.}

\emph{Thin support at large N.} Against a target of five clips per bin, N\ensuremath{\approx}1102 has 7, N\ensuremath{\approx}2963 has
2, N\ensuremath{\approx}6559 has 2 and N\ensuremath{\approx}13506 has 1; the survivor-N distribution is weighted toward small
graphs (median N=333, with 24 of 31 clips below N=1000). The headline fits rest on 4 dense
and 6 sparse points.

\emph{A guard violation, mechanism identified.} Two small-N clips (Assault036, N=97; Abuse037,
N=188) took a scene-sparse shortcut on 20\% and 10\% of queries, violating a pre-registered
guard that all timed queries traverse the PPR path. The mechanism is the margin test of §3:
when the top two shortlisted scenes are separated by more than \ensuremath{\tau}=0.015, or only one scene is
shortlisted, retrieval returns the anchor scene's exact top-k and skips subgraph PPR --- both
more likely as scene count falls, consistent with the two smallest clips being the affected
ones. The
shortcut is exact rather than approximate, so no retrieved result is degraded; but it is
also faster, so the affected queries bias the scene-sparse arm favourably. Both clips fall
outside the headline range, and refitting the full range without them raises the exponent
from 0.497 to 0.503 --- under 2\%, in the direction unfavourable to us.

\emph{Salience-weight provenance.} The corpus was built throughout at salience weights
(\texttt{luma\_diff\_weight}, \texttt{motion\_weight}, \texttt{luma\_entropy\_weight}) = (0.5, 0.3, 0.2), verified
per-clip from cached configuration snapshots. Despite its name, \texttt{luma\_diff\_weight} weights a
codec \textbf{packet-size residual} (\texttt{frame\_features{[}"packet\_size"{]}}), not a luma-difference
quantity; a genuine \texttt{luma\_diff\_energy} field exists on frame records and is diagnostic only,
never consumed by the scorer. These are the shipped defaults, so survivor-N values here are
comparable to retention figures computed under the production configuration.

\emph{Degree is flat against N, for this fit's edge mode only.} The exponent depends on per-scene
subgraphs not growing with N. Checked against all cached indices, degree does not drift
upward from N=24 to N=13,506 (r(N, deg\_mean) = \ensuremath{-}0.155), because scene count grows with N
(r=0.975) while scene size does not (r=0.135). This holds for the complete-block modes this
corpus is built under. It is \textbf{not} established at scale for the tiered mode our MLVU
results use: no cached index in this corpus was built under it, and the single tiered cache
measured since (one clip, N=613) confirms bounded degree there without establishing
degree-versus-N behaviour (§3).

\subsection{Tractability divergence}

At the largest sizes we measured, the dense arm does not return at all --- a capability
difference the exponent separation understates. Under a uniform 3600 s wall-clock and 29 GB
memory watchdog, the dense arm was censored at N=6,559 (Arrest047) and N=13,506 (Arson019),
with resident memory at 17.0 GB and 26.9 GB and still climbing at the cap; the scene-sparse
arm completed the same clips in 0.0104 s and 0.0210 s. Unlike the latency results
below, this is not subject to Amdahl dilution: a query that never returns cannot be rescued
by a fast downstream stage.

The censoring bounds what can be said. This is a tractability boundary, not a speedup ratio:
any ratio computed against a censored measurement is a lower bound on an unknown quantity.
Substituting the cap as a floor and refitting gives k \ensuremath{\geq} 2.597 for the dense arm --- a bound,
with no confidence interval attached.

\subsection{End-to-end accounting: where the speedup goes}

Our own headline retrieval ratio buys nothing at the sizes we tested. We report it because
the literature reports retrieval-side speedups routinely and end-to-end accountings rarely.

Retrieval-\textbf{mechanics} measurements at N=4,892 show a 1,104\ensuremath{\times} ratio between arms (7.9448 s
dense against 0.0072 s scene-sparse) under synthetic sampled query embeddings. That
ratio does not reach the user, and it is a retrieval-mechanics figure wherever it appears.

We measured the full query path --- query text to final answer --- at the same N with
\texttt{granite4:micro} (\textasciitilde3.4B, Q4\_K\_M) on CPU, \texttt{l2\_retrieve\_top\_k=30}, five real CLIP-text queries
per arm plus a discarded warm-up, and the caption cache reset before each arm. The primary
run is instrumented on a clean checkout (pinned commit, zero tracked-dirty files):

\begin{table}[t]
\footnotesize
\setlength{\tabcolsep}{4pt}
\caption{End-to-end query path at $N=4{,}892$ (per-stage medians).}
\begin{tabular}{lrr}
\toprule
component (median, s) & dense & scene-sparse \\
\midrule
retrieval mechanics & 8.535 & 0.030 \\
answerer (LLM) & 2.830 & 2.817 \\
captioning & 37.269 & 62.132 \\
L1 retrieval & 0.000 & 0.000 \\
Cerberus verification & 6.408 & 4.418 \\
\textbf{total} & \textbf{55.64} & \textbf{69.40} \\
\bottomrule
\end{tabular}
\end{table}

Stage medians are taken independently, so they need not sum to the totals. \textbf{End-to-end
ratio: 0.802\ensuremath{\times} --- the sparse arm is slower at the median.} Per-query caption cost spans
2.6--69.3 s depending on cache misses, so at five queries per arm these medians are
directionally indicative rather than precise. An earlier measurement under a different
captioner gives 0.778\ensuremath{\times}; because captioning is 85--93\% of the dominant stage, the two runs
corroborate \textbf{direction only, not magnitude}.

Captioning is 85\% of the dominant stage in the dense arm and 93\% in the sparse arm. L1
retrieval costs exactly zero on every query in both arms --- it populates an in-memory
dictionary over already-retrieved frames --- so the stage is in practice a two-way split
between captioning and verification. Within it, frame access is not the cost: on one
cold-cache query at the same N (30 misses, 272 frames decoded, 27 distinct GOPs), the decode
loop cost 0.815 s against 114.6 s of captioning, under 1\% of the stage.

The retrieval ratio in this run is approximately 291\ensuremath{\times}, against the 1,104\ensuremath{\times} above. The two are
measured under different query types --- real CLIP text encoding here, synthetic embeddings
there --- and must not be conflated.

One observation matters more than the headline ratio: \textbf{the LLM is not the bottleneck.} The
answerer costs 2.83 s and is identical across arms, as expected when both feed the same
top-k context to the same model. The dominant term is captioning, which is O(top\_k) rather
than O(N) and so does not shrink as the graph sparsifies. Sparsifying makes the retrieval
step negligible while leaving the largest term untouched.

\textbf{Crossover.} The retrieval-only crossover --- the size past which scene-sparse retrieval
mechanics cost less than dense --- sits at N\ensuremath{\approx}24 from the §5.1 fits. We do not project an
end-to-end crossover. Doing so requires assuming the constant stage cost is equal across
arms, which this run's own data contradict (§5.4), and a projection built on that assumption
returns \textasciitilde1.15\ensuremath{\times} at N=4,892 where we measure 0.802\ensuremath{\times}. The tractability boundary of §5.2 is the
stronger statement and we rest on it.

\subsection{Caption-cache locality (negative result)}

The sparse arm's captioning costs more than the dense arm's at identical top\_k. Over 20 real
text queries per arm at N=4,892, with top\_k verified byte-identical across arms (30 frames
on all 40 queries) and an identical query set, the paired clip-level bootstrap difference is
\textbf{+9.30 s per query, 95\% CI {[}3.52, 16.22{]}} (seed 42, B=10,000), excluding zero.

\begin{table}[t]
\footnotesize
\setlength{\tabcolsep}{4pt}
\caption{Caption-stage cost per query (20 queries per arm).}
\begin{tabular}{lrr}
\toprule
(mean per query) & dense & scene-sparse \\
\midrule
caption cache misses & 3.85 & 7.75 \\
frames decoded & 38.3 & 73.15 \\
per-frame caption wall time & 2.535 s & 2.520 s \\
verify calls & 1.05 & 1.00 \\
\bottomrule
\end{tabular}
\end{table}

Per-frame cost is equal across arms; the entire difference is miss count. We tested and
\textbf{rejected} the hypothesis that the sparse arm retrieves more temporally dispersed frames:
sparse retrievals touch marginally fewer distinct scenes (23.8 against 26.8), have a smaller
mean pairwise frame-index distance (2,732 against 3,070), and dispersion correlates
essentially not at all with caption time (pooled r = \ensuremath{-}0.058, n=40).

The mechanism is cross-query cache locality. The dense arm re-retrieves an overlapping pool
of high-salience frames largely independent of the query --- its miss count reaches zero by
the sixth distinct question --- whereas the sparse arm returns query-specific frames, so
successive questions touch largely disjoint frame sets and the cache warms slowly.

The double edge is worth stating: the dense arm's caching advantage is a consequence of its
retrieval being \emph{less responsive to the query}. The same property that makes it cheap to
cache makes it a worse retriever. That does not make the cost imaginary --- the caption cache
persists across a session in production --- so the honest statement is that the sparse graph
retrieves far more cheaply \emph{and} benefits less from cross-query caption reuse. A build-time
caption prefetch over high-centrality nodes would plausibly eliminate it, given a 0.212 s
construction budget; we have not implemented or measured that.

\section{Correctness Floor}

This section does not claim competitive accuracy. It claims that the pipeline, with no
generative model in its construction path and a \textasciitilde3.4B answerer on CPU, answers within the
weakly-supervised band. It does not measure the 0.212 s construction of §4: §6.1 runs a flat
dense graph and §6.2 the tiered one (§3.5). The link to §4 is the identity gate of §4.1,
which shows the cheap build retrieves identically to dense-then-prune; no answer accuracy is
measured on either (§8).

\subsection{Grounded QA on NExT-GQA}

\textbf{Split provenance.} Our pool is drawn entirely from NExT-GQA's \emph{validation} grounding
subset (\texttt{gsub\_val}) intersected with our cached videos: 526 questions over 86 videos,
partitioned at video level into a 59-video tuning half (406 questions) and a \textbf{27-video
held-out validation half} (120 questions). We call the latter a held-out \emph{val} split, not a
test split: the official test set (5,553 questions over 990 videos) is untouched by any
experiment in this paper.

On that split, \textbf{Acc@GQA is 0.1667, 95\% CI {[}0.088, 0.243{]}}, with Acc@QA 0.375 {[}0.275, 0.477{]}; the answerer is \texttt{granite4:micro} at temperature 0 on CPU, and the option parser
succeeded on 120/120 questions. The result decomposes cleanly: 42 of 120 questions are
correctly grounded, 45 answered correctly, 20 both --- 0.350 \ensuremath{\times} 0.476 = 0.167.

\textbf{Grounding and correctness are coupled.} P(correct \textbar{} grounded) = 0.476 against
P(correct \textbar{} ungrounded) = 0.321, a \textasciitilde15-point gap that replicated across two independent
samples even though absolute levels fell between them. What replicates is the gap, not the
levels.

\textbf{But the two factors are readouts of one shared retrieval event, not independent
measurements.} \texttt{packet\_size} enters that event's single PageRank call twice: through the
personalization vector at weight (1\ensuremath{-}\ensuremath{\lambda})=0.5, and through the edge weights at the \ensuremath{\beta} slot under
the default motion mode. Grounding and answering then read the identical retrieved-frame
list. This does not undermine the gap: ``grounded'' means the correct region's frames were
present in the answerer's context, a real causal channel however retrieval found them. What
it narrows is the stronger reading of grounded as \emph{retrieved for a query-semantic reason} ---
since up to half the seed is a codec artifact uncorrelated with the question text, a question
can be grounded because \texttt{packet\_size} ranked the right region highly. This bears on §2.2's
framing of grounding as the guard against being right for the wrong reasons: the guard
establishes that the right evidence was present, not that it was found for the right reason.
A residual confound, in which some codec characteristic correlated with favourable retrieval
also correlates with question difficulty, cannot be ruled out from a code trace; no control
analysis for it exists, and we disclose that as unaddressed.

\textbf{Neither stage is negligible.} Perfect grounding would raise Acc@GQA to 0.476 (+0.31); a
perfect answerer over current grounding would give 0.350 (+0.18).

\textbf{Against baselines.} Our held-out advantage over uniform frame sampling in answer accuracy
is +0.120, 95\% CI {[}0.000, 0.248{]} --- the interval touches zero. Against \emph{random} sampling it is
+0.127 {[}+0.036, +0.229{]}, which separates. The honest statement: retrieval beats random frame
sampling; against uniform sampling the advantage is positive but not statistically separated
at n=120.

\begin{table}[t]
\footnotesize
\setlength{\tabcolsep}{4pt}
\caption{NExT-GQA placement. Published rows use the full 5,553-question test set; ours is a 120-question held-out validation split.}
\begin{tabular}{llrrr}
\toprule
method & params & Acc@GQA & mIoP & IoP@0.5 \\
\midrule
Temp{[}CLIP{]} NG+ & 130M & 16.0 & 25.7 & 25.5 \\
SeViLA* & 4B & 16.6 & 29.5 & 22.9 \\
LangRepo & 12B & 17.1 & 31.3 & 28.7 \\
FrozenBiLM NG+ & 1B & 17.5 & 24.2 & 23.7 \\
VideoMind-2B & 2B & 25.2 & 36.4 & 32.6 \\
MUPA-2B & 2B & 28.7 & 39.1 & 38.7 \\
MUPA-7B & 7B & 30.3 & 41.4 & 39.4 \\
\textbf{IRIS (ours)} & \textasciitilde3.4B, CPU & \textbf{16.7} & --- & --- \\
\bottomrule
\end{tabular}
\end{table}

\textbf{This is indicative, not a leaderboard entry, and the splits differ.} Ours is n=120 from a
held-out validation split with an \textasciitilde8-point interval; every published row is the full
5,553-question \emph{test} set. We place the row to situate the system, not to rank it, and claim
to have beaten nothing in the table. A comparable number would require the official test
split, which remains available to us precisely because we never touched it. The held-out val
split, by contrast, is \textbf{burned}: both permitted touches are used, so no further measurement
on those 27 videos is possible without a new split --- a constraint that binds the
answerer work this section identifies as the live target.

We omit mIoP and IoP@0.5 from our row deliberately. They are measured and
convention-checked, but this is a correctness floor rather than a grounding claim, and
printing a grounding magnitude from 120 validation questions beside full-test figures would
invite exactly the comparison we decline. Our Acc@GQA is keyed to IoP@0.5, which is identical (0.3202 on
the tuning half) under the benchmark's max-per-span convention and the union convention an
earlier version of our evaluation used, so the reported figure is convention-invariant.

\subsection{Multi-task long-video QA on MLVU}

We also evaluate MLVU's multiple-choice tasks, restricted to the six third-person ones: the
egocentric task is excluded because our compressed-domain signals assume a largely static
camera, and the generation tasks fall outside our multiple-choice harness. 150 questions, 25
per task, videos capped at 600 s, codec segmentation, scene-sparse graph.

\textbf{M-Avg 0.340}, zero ingest failures across 145 unique videos, overall option-parse failure
3.3\%.

\begin{table}[t]
\footnotesize
\setlength{\tabcolsep}{4pt}
\caption{MLVU per-task accuracy (150 questions, 25 per task).}
\begin{tabular}{lrrr}
\toprule
task & accuracy & chance & parse-fail \\
\midrule
Topic Reasoning & 0.720 & 0.250 & 4.0\% \\
Plot QA & 0.400 & 0.250 & 0.0\% \\
Anomaly Recognition & 0.320 & 0.250 & 0.0\% \\
Needle QA & 0.280 & 0.250 & 12.0\% \\
Action Order & 0.160 & 0.250 & 4.0\% \\
Action Count & 0.160 & 0.250 & 0.0\% \\
\bottomrule
\end{tabular}
\end{table}

\textbf{Two tasks fall below chance, and the cause is the answerer rather than retrieval.} Action
Order and Action Count sit at 0.160 against a 0.250 floor, with parse failure at 4\% and 0\%,
so the model is confidently wrong rather than unparsed. Both require what top-k retrieval
does not supply: chronological ordering needs reasoning over relative time, and exhaustive
counting needs complete coverage rather than the most relevant evidence. This is the same
failure mode as NExT-GQA's directional-temporal questions --- a limitation of the answerer at
this scale rather than of the graph (§8). Needle QA's 12\% parse-failure rate is the highest of
the six and its accuracy sits closest to chance among tasks that clear it; the two may be
related, and we claim no meaningful margin there.

\subsection{What this section supports}

The system answers, in the band occupied by weakly-supervised methods with comparable or
larger models, on CPU with no training and no generative model in construction --- though not
yet on the 0.212 s configuration of §4 (§8).
Where it fails --- directional temporal reasoning, exhaustive counting --- the failure is in the
answerer rather than the representation.

\section{What Does Not Matter}

The results in this section are negative. We report them because each was tested against a
stated kill criterion, each criterion triggered, and together they delimit our own
contribution: they are why §4 claims \emph{cheap} construction rather than \emph{better} construction.
The two criteria differ in standing. §7.1's was pre-registered. §7.2's was recorded, but it
entered version control alongside the run it governs, so its precedence cannot be verified
and we do not describe it as pre-registered.

\subsection{Segmentation placement, once the budget is fixed}

Our scene boundaries derive from compressed-domain residual peaks (§3.2). Does that
content-adaptive placement produce a better graph than a content-blind rule? We compared four
boundary strategies with retrieval, answerer, question set and seed held constant: \textbf{codec}
(residual peaks), \textbf{equal-count} and \textbf{equal-time} (both at the segment count codec produced
for that video, placed for equal survivor counts and equal intervals respectively), and
\textbf{fixed-60s}, the interval used by comparable pipelines. The matched-count arms isolate
\emph{where} boundaries fall from \emph{how many} there are.

On short videos (\ensuremath{\leq}600 s; six MLVU tasks, 150 questions, 145 videos) codec did not lead: M-Avg
0.340 codec, 0.393 equal-count, 0.320 equal-time, 0.353 fixed-60s --- the simplest matched-count
control 5.3 points ahead. Short videos yield few segments, so placement has little room to
matter; a long-video run (600--1800 s, AR and PQA, 45 videos) tests the regime where it should,
and the premise holds: codec produced a median of 581 scenes per video (range 174--1319)
against single digits on short videos.

\begin{table}[t]
\footnotesize
\setlength{\tabcolsep}{4pt}
\caption{Long-video segmentation ablation (600--1800\,s).}
\begin{tabular}{lrrr}
\toprule
arm & AR (n=17) & PQA (n=91) & M-Avg \\
\midrule
codec & 0.412 & 0.451 & \textbf{0.431} \\
equal-count & 0.235 & 0.451 & \textbf{0.343} \\
\bottomrule
\end{tabular}
\end{table}

The codec-minus-equal-count difference is \textbf{+0.088 M-Avg, 95\% CI {[}\ensuremath{-}0.009, +0.204{]}}
(video-clustered bootstrap, seed 42, B=10,000). The interval includes zero, the criterion
required it to exclude zero, and a tie was registered as failure. The difference did move in
the predicted direction (\ensuremath{-}0.053 on short videos to +0.088 on long), but all of it is carried
by anomaly recognition (+0.176, CI {[}0.000, 0.353{]}, n=17) while plot QA is exactly flat
(+0.000, CI {[}\ensuremath{-}0.101, +0.102{]}). \textbf{We therefore do not claim that codec-derived boundaries
produce a better graph.} A domain-concentrated benefit remains possible, untested at a sample
size that could resolve it.

The finding a practitioner should take: once the segment budget is fixed, boundary placement
is within noise across four strategies spanning content-adaptive to content-blind. That
licenses using the cheapest available segmentation, and it is a stronger argument for the
pipeline than a codec win would have been, because it does not depend on our signal being
special.

\subsection{Frame admission versus uniform sampling}

Do codec signals help \emph{select} which frames enter the graph? Three tests on the 19 annotated
UCF-Crime anomaly videos say they do not.

\textbf{Coverage metrics are uninformative without a matched control.} Uniform admission saturates
gold-window coverage (M1 = 1.0000) at every budget swept, down to 5\% retention, so any
selector appears to cover 99\%+ of gold windows: the metric measures annotated-window length,
not selection quality. We flag this because we previously reported such a figure ourselves.

\textbf{At matched budget, codec admission is indistinguishable from uniform.} At 10.5\% retention
the paired differences on the two informative metrics are +0.256 pp (95\% CI {[}\ensuremath{-}0.236, +0.710{]})
and +0.780 pp ({[}\ensuremath{-}0.005, +1.753{]}), both spanning zero; the admitted set sits at the
label-blind identity, mean(M2 \ensuremath{-} retention) = +0.0017 (sd 0.0099).

\textbf{Whether that tie is an artifact of a wasted budget is unresolved.} The score might clump
its picks onto near-duplicate frames, in which case the tie would reflect poor temporal
coverage rather than an uninformative signal. A post-hoc diagnostic --- explicitly not
pre-registered --- measured dispersion at a matched budget of k=503: the codec top-k arm sampled
42 distinct temporal locations (8.4\% of budget, mean run length 13.19) against uniform's 503
(100\%, mean run 1.00), because the action score is hold-forward propagated from the retained
tier and is therefore a step function of roughly 9.5 frames per plateau, so top-k degenerates
into taking whole plateaus from their earliest frame. \textbf{This does not dispose of the
objection}: the diagnostic cannot separate an uninformative signal from a budget spent on
near-duplicates, and the plateau-dedup contrast that would is specified but not run.

\textbf{Shot geometry adds nothing either.} Segmenting each video into shots directly from the
packet curve (zero decode) and sampling one frame per shot at its midpoint gives, against
uniform, +0.001 pp ({[}\ensuremath{-}0.416, +0.421{]}) and +0.141 pp ({[}\ensuremath{-}0.608, +0.972{]}) --- half-widths that
\emph{exclude} a 1 pp effect, making this the one properly powered contrast in the group and a
genuine null. The corresponding shot-plus-score contrast was underpowered (half-widths 1.69
and 2.79 pp) and we report it as such rather than as a tie.

\textbf{Power.} The annotated corpus is 19 videos, so each is 5.3 pp of any aggregate and the
resolution floor is roughly 1 pp; only more labelled anomaly footage changes that, and we did
not re-run in search of significance.

\subsection{What these negatives buy}

Neither \emph{where} we cut nor \emph{which frames we keep} is doing the work. What does is structural:
the graph is block-diagonal, so it is cheap to build and cheap to query, and that property is
independent of how the blocks are chosen --- narrower than the claim we set out to make, and
more portable. A pipeline adopting the construction needs the block structure, not our codec
signal, our segmentation, or our admission policy.

\section{Limitations}

Negative results appear inline where they arose (§5.3, §5.4, §7). Four of the constraints
that bound our claims are also stated there and only indexed here: the retrieval speedup is
not an end-to-end speedup (0.802\ensuremath{\times}, §5.3); sparse retrieval pays a caption-cache penalty of
+9.30 s per query (§5.4); the scaling fit rests on 4 dense and 6 sparse points at large N
(§5.1); and accuracy sits in the weakly-supervised band, 8--14 points below current agentic
methods (§6.1). The construction result carries no equivalent end-to-end caveat.

\textbf{The construction saving is not yet measured on the evaluated configuration.} This is the
most substantive open item in the paper. No accuracy number we report is measured on the
complete-block configuration: NExT-GQA uses a flat dense graph, to which block-diagonal
construction cannot apply, and MLVU the tiered configuration, whose neighbour selection scans
all pairs regardless of how sparse its output is (§3.5, §4.3). Bounding those scans to the
scene partition would plausibly transfer the saving to the tiered configuration, but it yields
a different graph, so it needs accuracy validation rather than an identity proof. Until then, §4 is a construction
result about a configuration we do not evaluate.

\textbf{Directional temporal reasoning fails, and not in the graph.} Action Order and Action Count
fall below chance with near-zero parse failure, and NExT-GQA's before/after questions show the
same pattern (§6.2). Both need relative temporal ordering or exhaustive coverage, which top-k
retrieval does not supply and which no change to graph construction addresses.

\textbf{Scene shortlisting bounds recall.} A scene excluded by the \(\lceil\sqrt{S}\rceil\)-width shortlist cannot
contribute evidence downstream (§3.4). Whether that width is recall-safe at large S is not
evaluated.

\textbf{Evaluation splits are constrained.} The NExT-GQA held-out \emph{validation} half is burned
(§6.1), which binds future answerer work more than the results reported here. MLVU is 150
questions across six tasks with videos capped at 600 s.

\textbf{Answerer reproducibility is scoped.} §6.1's answerer ran on llama-server at temperature 0.
A byte-identical re-query check (639/639) passed for that serving configuration, but on a
different build and machine, so it corroborates rather than establishes §6.1's determinism.
The MLVU results (§6.2, §7.1) ran through Ollama, whose client never sent the prompt-cache
disable flag and has no equivalent to single-slot serving, and every MLVU question was
answered exactly once, so \textbf{their reproducibility under re-query is undetermined}. §5.3--§5.4
also ran through Ollama, but report only timings and counts, which this does not affect.

\textbf{Encoding profiles are uncharacterised.} UCF-Crime is redistributed web video, already
compressed before it reached us, so every codec-derived signal we read is a second-generation
measurement of an encoding history we neither observed nor control; part of the corpus is
further x264-transcoded, and the flagship VIRAT efficiency clip is \texttt{mpeg4} where the accuracy
corpus is H.264. The salience channel (§3.1), the codec-derived scene boundaries behind §4,
and both negatives in §7 are conditioned on the profiles this corpus happens to carry. This
is a scope statement, not a weakness we have evidence of.

\textbf{Hardware and provenance.} The CPU-only claim rests on a \texttt{torch.cuda.is\_available()} flag
captured from a CPU-only torch wheel, plus circumstantial path evidence that ingest ran on a
Windows machine rather than the project's Linux GPU box; it establishes that no CUDA device
was used, not that none was present. The zero-forward-pass claim (§3.1) is a property of the
code path and is unaffected. The \textasciitilde6.8\ensuremath{\times} peak-memory ratio is not tied to a commit, unlike the
wall-clock ratio, and the build-cost harness was committed after its run (§4.2).

\section{Conclusion}

We asked whether a long-video retrieval graph must be expensive to build. It need not be.
Block-diagonal construction produces exactly the graph that dense construction produces and
then prunes --- same nodes, edges, weights at zero tolerance, and PageRank --- while visiting
only the pairs it keeps, making construction \textasciitilde234\ensuremath{\times} faster and \textasciitilde6.8\ensuremath{\times} smaller in memory on CPU,
with no generative model in the construction path. The same block structure separates query
latency into quadratic and sublinear regimes, and at large sizes into a boundary where dense
retrieval does not complete.

We have been deliberate about what this does not establish: the retrieval speedup does not
reach the user because captioning dominates the query path, sparse retrieval pays a
caption-cache penalty, and accuracy sits in the weakly-supervised band. Our two falsified
hypotheses --- codec-derived boundaries (pre-registered) and codec-based admission (criterion
recorded, registration date unverifiable) --- show that neither where we cut nor which frames
we keep does the work. \textbf{The block structure alone does}, which is what makes the result
portable. For retrieval over long video, the intermediate graph does not have to be
generated, and does not have to be expensive. It has to be structured.

\begin{acks}
% GenAI disclosure (ACM policy) -- wording to be agreed by all authors before submission.
\end{acks}
% REFERENCES: bibliography not yet built; inline citations are placeholders.
\end{document}
